%% ===================================================================
%% Chapter 2. Related Works
%% ===================================================================

\chapter{Related Works}
\label{ch:related}

This chapter reviews the three research areas on which this thesis builds: autonomous drone interception, competitive multi-agent reinforcement learning with self-play, and online adaptation of learned policies.

\section{Autonomous Drone Interception}

\subsection{Classical Guidance and Capture Systems}
\noindent
Counter-UAS systems span detection, tracking, and neutralization, with physical capture by an interceptor drone being one of the few options that is safe to use in populated areas \cite{chauhanNationsDefenseComprehensive2025}. Vrba et al.\ \cite{vrbaAutonomousCaptureAgile2022} built a complete LiDAR-based system that predicts the target trajectory and repositions a net-carrying UAV to catch it, achieving top results in the MBZIRC 2020 challenge. Pliska et al.\ \cite{pliskaSafeMidAirDrone2024} proposed Fast Response Proportional Navigation together with an interacting-multiple-model filter for target estimation, and showed that estimation quality is as important as the guidance law itself. Model-based trajectory generation for fast flight \cite{zhouRobustEfficientQuadrotor2019} and decision-level allocation of interceptors to swarm threats \cite{palmasReinforcementLearningDecisionLevel2025} complement these approaches. A common limitation is that the target is treated as a process to be estimated and extrapolated rather than as an agent that reacts to the interceptor.

\subsection{Reinforcement Learning for Agile Flight}
\noindent
Deep RL has enabled quadrotor control at the limits of the platform. Policies trained in simulation have performed acrobatic maneuvers \cite{kaufmannDeepDroneAcrobatics2021} and defeated human world champions in drone racing \cite{kaufmannChampionlevelDroneRacing2023a}, with later systems winning competitions using only monocular vision \cite{bahnamMonoRaceWinningChampionLevel2026} or operating in uninstrumented arenas \cite{boselloYourOwnProLevel2026}. Kaufmann et al.\ \cite{kaufmannBenchmarkComparisonLearned2022} showed that policies commanding collective thrust and body rates (CTBR) transfer to hardware more robustly than those commanding individual rotor thrusts, and Chen et al.\ \cite{chenWhatMattersLearning2025a} identified practical factors---including action smoothness regularization, system identification, and large batch sizes---that enable zero-shot sim-to-real transfer on small quadrotors such as the Crazyflie. GPU-parallel simulators, notably OmniDrones on Isaac Sim \cite{xuOmniDronesEfficientFlexible2023}, have made training at this scale practical. Related lines of work address safe navigation among dynamic obstacles \cite{xuNavRLLearningSafe2025, ugurluSimtoRealDeepReinforcement2022}, differentiable-physics training \cite{zhangLearningVisionbasedAgile2025, gurumurthyDeepOffPolicyIterative2023}, and generalization across tasks \cite{greenBridgingPerformanceGeneralization2026}.

\subsection{Reinforcement Learning for Interception}
\noindent
RL-based interception has been studied for missiles using line-of-sight observations \cite{huInterceptGuidanceManeuvering2023}, for fixed-wing UAVs in controlled airspace \cite{giralInterceptingUnauthorizedAerial2024}, and for multirotors using visual input \cite{darwishDroneNavigationTarget2023} or obstacle-aware planning in urban scenes \cite{lucotteInterceptionAdversarialDrones2026}. Multi-target and multi-agent variants combine RL with fuzzy logic \cite{xiaUAVMultiDynamicTarget2024}, handle switching communication topologies \cite{fanSwitchingawareMultiagentDeep2023}, or coordinate several pursuers with target-prediction networks \cite{zhangGameDronesMultiUAV2023, chenOnlinePlanningMultiUAV2025}. Control-theoretic RL formulations of the pursuit--evasion differential game add safety and robustness guarantees \cite{wangReinforcementLearningPursuitevasion2025}. In most of these works the target follows scripted or randomized motion, so the learned pursuer is not exposed to an evader that optimizes against it.

\section{Competitive Multi-Agent Reinforcement Learning and Self-Play}

\subsection{Non-Stationarity in Competitive Learning}
\noindent
When both sides of a game learn, each agent faces a non-stationary environment, which breaks the convergence assumptions of single-agent RL \cite{zhangMultiAgentReinforcementLearning2021, liangReviewMultiAgentReinforcement2025}. Centralized-critic methods such as MADDPG \cite{loweMultiAgentActorCriticMixed2017} and multi-agent PPO \cite{yuSurprisingEffectivenessPPO} mitigate this in part, while opponent-modeling approaches explicitly predict how the opponent will behave \cite{tianMultiagentActorcriticTime2023}. In purely competitive settings, naive self-play against only the latest opponent is prone to overfitting, forgetting, and cyclic strategies; benchmarks such as FightLadder \cite{liFightLadderBenchmarkCompetitive2024} show that producing non-exploitable agents remains difficult. Feng et al.\ \cite{fengMultiagentEmbodiedAI2026} highlight the same challenges for embodied multi-agent systems more broadly.

\subsection{Population-Based Self-Play}
\noindent
Population-based training keeps a pool of previous and diverse policies and trains new policies against mixtures of that pool. AlphaStar \cite{vinyalsGrandmasterLevelStarCraft2019} used a league of main agents and exploiters with prioritized fictitious self-play (PFSP), which samples opponents in proportion to how often they beat the learner. In air combat, Chen et al.\ \cite{chenAutonomousDogfightDecisionMaking2025} proposed automatic opponent sampling with a simulated-annealing-inspired curriculum and Boltzmann meta-solver, and Kong et al.\ \cite{kongUAVAutonomousAerial2020} addressed robustness to observation errors.

For quadrotor pursuit--evasion, Xiao and Feroskhan \cite{xiaoLearningMultipursuitEvasion2024} trained pursuers and evader asynchronously over multiple stages, and Roncero et al.\ \cite{ronceroLearnedControllersAgile2026} extended this to a PSRO-style population with Hedge sampling over win rates, training one side per generation against frozen opponents. AgilePE \cite{tangAgilePEAutonomousUAV2026a} and Gavin et al.\ \cite{gavinAgileInterceptionFlying2026, gavinInterceptingAgileTarget2026} used PFSP with CTBR control and demonstrated zero-shot transfer to real quadrotors, while Geles et al.\ \cite{gelesSuperhumanSafeAgile2026} applied league-based self-play to multi-drone racing. Xu et al.\ \cite{xuDRLFrameworkAutonomous2025} combined self-play PPO with curriculum learning across spacecraft and drone domains. These works establish population-based self-play as an effective tool for aerial pursuit--evasion, but they rank opponents almost exclusively by win rate.

\section{Online Adaptation of Learned Policies}
\noindent
A policy trained in simulation is typically deployed with frozen weights. To handle conditions not seen in training, recent work adapts the policy online. Pan et al.\ \cite{panLearningFlyRapid2026} refine a residual dynamics model from real data and update the policy through differentiable simulation within seconds. Xing et al.\ \cite{xingContinualRobotPolicy2026} infer a latent dynamics condition from recent interaction with a recurrent encoder and condition the policy on it. Pertin and Chang \cite{pertinNSTRLResidualReinforcement2026} learn residual corrections on top of an approximately optimal regulator. These methods target changes in the robot's own dynamics and rely on gradient updates or learned encoders.

Reservoir computing offers a cheaper route. An echo state network (ESN) uses a fixed, randomly initialized recurrent reservoir whose state encodes the input history; only a linear readout is learned. Yoshimura and Tanaka \cite{yoshimuraOnlineAdaptiveReinforcement2026} attached an ESN to an RL policy and updated its readout online by recursive least squares (RLS), achieving adaptation to abrupt changes within a few control steps without backpropagation, pretraining, or privileged information. Their evaluation considered non-stationary dynamics in single-agent benchmarks; applying the same mechanism to a non-stationary \emph{opponent} in an adversarial aerial task has not been explored.

\section{Positioning of This Thesis}
\noindent
This thesis combines the two threads above. Like \cite{ronceroLearnedControllersAgile2026, tangAgilePEAutonomousUAV2026a, gavinInterceptingAgileTarget2026}, it trains CTBR quadrotor policies by population-based self-play, but it replaces win-rate-only opponent sampling with a multi-metric hardness score that also accounts for capture time, control effort, and the learner's own failure rate. Building on \cite{yoshimuraOnlineAdaptiveReinforcement2026}, it equips the pursuer with an ESN whose RLS readout predicts the evader's relative motion online, providing a gradient-free mechanism for adapting to the specific evader encountered at test time.
